% 本文件由 LaTeX 排版 Agent 根据有效 translation.json 自动生成。
% author=Augustine of Hippo; work_id=1313; title=致康森提乌斯，驳说谎

\chapter{致康森提乌斯，驳说谎}

% structure: p0001 source=h1 target=omitted reason=page-title-already-rendered-by-parent

\EditorialNote{出自《订正录》卷二，第六十章：此后我又写了一部《驳说谎》的书，其缘起如下。为了发现那些认为不仅可以用否认和说谎，甚至可以用伪誓来隐藏其异端的普里西利安派异端，有些天主教徒认为，他们自己也应当假装成普里西利安派，以便潜入其巢穴。为禁止这种行为，我撰写了本书。它以\Quote{Multa mihi legenda misisti}开头。}\par

1. 至亲的弟兄康森提乌斯，你寄来许多要我阅读的文字，实在很多。在我准备答复之时，先被一件又一件更紧急的事务所打断，一年已经过去，并把我逼入这样的困境，以致我必须尽我所能地答复，免得如今航行季节正好，送信者渴望返回，我却将他久留。因此，当莱昂纳斯，天主的仆人，从你那里带给我的所有东西，在收到后不久，以及后来在我口述此答复时，我都展开并通读，并尽我所能地仔细权衡，我极其喜爱你的口才、对圣经的记忆、才智的敏锐，以及你谴责疏忽职守的天主教徒时的愤慨，和你在面对潜伏的异端时咬牙切齿的热忱。但我并不相信，借由我们撒谎来把他们从藏身之处挖出来是正当的。因为我们费尽心机追踪他们、追捕他们，最终目的是什么？不就是把他们抓住并带到光天化日之下，要么教导他们真理，要么至少用真理定他们的罪，不让他们伤害他人吗？因此，其目的是使他们的谎言被清除或避开，并使天主的真理增长。那么，我如何能借谎言正直地追讨谎言呢？难道是用抢劫来追讨抢劫，用亵渎来追讨亵渎，用奸淫来追讨奸淫吗？\BibleQuote{但若天主的真理因我的谎言而更加彰显}，我们是否也要说：\BibleQuote{让我们作恶以成善}呢？你看宗徒是如何憎恶此事的。因为\Quote{我们撒谎，以便将异端撒谎者引向真理}，这不就是\BibleQuote{让我们作恶以成善}吗？或者，谎言有时是好的，或者有时谎言不是邪恶？那么，为何经上写着：\BibleQuote{主，你恼恨一切作恶的人；你必消灭一切说谎的人}？因为他没有排除某些人，也没有笼统地说：\BibleQuote{你必消灭说谎的人}，以致允许我们理解某些人，而非所有人；而是他宣告了一个普世的判决，说：\BibleQuote{你必消灭一切说谎的人}。或者，因为经上没有说\BibleQuote{你必消灭一切说谎的人}中的\Quote{一切}，或者说\BibleQuote{你必消灭任何说谎的人}，所以就要认为有些谎言是被允许的吗？也就是说，有些谎言和说这些谎言的人，天主不会消灭，而是消灭所有说不义谎言的人，而不是任何谎言，因为还有一种正义的谎言，它应当值得称赞，而非罪行？\par

2. 难道你没有察觉，这种推理大大助长了那些我们正大费周章用谎言去捕捉的猎物本身吗？因为，正如你已表明的，这正是普里西利安派的观点；他们为了证明这一点，引用圣经中的见证，鼓励其追随者撒谎，仿佛以圣祖、先知、宗徒、天使为例；甚至毫不犹豫地加上主基督自己；并认为除非他们宣称真理是撒谎者，否则就无法证明他们的虚假是真实的。必须驳斥这一点，而不是模仿；我们也不应在这一恶事上与普里西利安派同伙，他们在其中被证实比其它异端更坏。因为他们独自，或者至少他们最大程度地，被发现将撒谎作为教条，以便隐藏他们所谓的真理：并且他们认为如此大的恶事是正义的，因为他们说，必须在心中持守真理，但对陌生人用口说出虚假之事并无罪过；并且经上写着：\BibleQuote{他心中讲真理}，仿佛这足以使人成义，即使一个人用口说出谎言，当听到的不是他的邻人而是陌生人时。为此，他们认为宗徒保禄在说了\BibleQuote{你们要弃绝谎言，各人与邻人说实话}之后，立刻紧接着说：\BibleQuote{因为我们是互相为肢体}。意思是，对于那些在真理的团契中不是我们的邻人，也不是我们的肢体的人，说谎是合法且正当的。\par

3. 这句判词侮辱了神圣的殉道者，更确切地说，完全取消了神圣的殉道。因为按照这些人的看法，他们若更公义、更明智地行事，就不该向迫害者宣认自己是基督徒，并因宣认而使他们成为凶手；反而该通过说谎和否认自己的身份，既保全肉身的方便和内心的意图，又不让那些人实现他们心中所筹划的恶行。因为他们并非基督信仰中的邻人，以至于有责任用口说出心里的话；反而正是一\Emph{真理}本身的仇敌。因为如果\Person{耶胡}（他们似乎精明地将他挑出来，视作说谎的榜样）假称自己是\Person{巴耳}的仆人，以便杀死\Person{巴耳}的仆人：那么按照他们的悖谬，在迫害时期，基督的仆人岂不更可公义地假称自己是魔鬼的仆人，以便魔鬼的仆人不要杀死基督的仆人；并祭献偶像使人不被杀，如果\Person{耶胡}祭献\Person{巴耳}是为了杀人？因为根据这些说谎者的杰出教义，如果他们在身体上虚假地佯装崇拜魔鬼，而心中保持对天主的崇拜，那对他们有何损害呢？但殉道者们并非如此理解宗徒，真正的、神圣的殉道者。他们看见并持守了所记载的：\BibleQuote{心里相信，可得以成义；口里宣认，可得以救恩；}\newline{}和\BibleQuote{在他们口中找不出谎言；}\newline{}于是他们无可指责地离去，到了那不再害怕被说谎者试探的地方，因为在天上的集会中，他们不再有说谎者，无论是外人还是邻人。至于那\Person{耶胡}，通过亵渎的谎言和亵渎的祭献来搜寻亵渎之人以便杀死他们，他们不会效仿他，即使圣经对他的品行毫无所言。但是，既然记载说他的心在天主面前不正，那么他在某些顺服上（关于彻底消灭\Person{阿哈布}家族，他出于自己统治的欲望而表现）所获得的暂时王国的短暂报酬，又有什么益处呢？但愿殉道者的真实判词是你所辩护的：我以此劝勉你，我的兄弟，使你反对说谎者，不是作为说谎的导师，而是作为真理的维护者。因为我恳求你，请留心听我所说的话，以便你发现那种被认为适合教导的东西是多么需要避免，它虽然对不敬虔之人怀着可称赞的热忱，为了他们能被捕获和纠正，或避免他们，但却是过于轻率地。\par

4. 谎言有多种，诚然，我们普遍应当憎恨一切谎言。因为没有谎言不与真理相悖。因为，如同光明与黑暗、虔诚与不虔、正义与不义、罪与善行、健康与软弱、生命与死亡，真理与谎言也是彼此对立的。因此，我们爱前者有多深，就应恨后者有多深。然而事实上，有些谎言相信并无害处：虽然用这种谎言去欺骗，对说谎者有害，对相信者却无害。例如，如果那位弟兄、天主的仆人\Person{弗龙托}，在他给你的报告中（但愿不会！）说了些假话，他肯定伤害了自己，而不是你，尽管你并无不义地相信了他所说的一切。因为，不论那些事是否真的发生，它们没有任何东西，如果一个人相信它如此，即使并非如此，他也不会因真理的准则和永恒救恩的教义而被判定有罪。相反，如果一个人说了一个谎言，任何人若相信它就会成为反对基督训导的异端者，那么说谎者造成的伤害越大，相信者就越悲惨。那么请看，这是何等的事，如果我们为了捕捉同一训导的仇敌，为了将他们引向真理而远离真理，竟要对着基督的训导说出一个谎言，任何人相信它都会丧亡；更糟的是，当我们用谎言捕捉说谎者时，我们实际上在教授更严重的谎言。因为，他们说谎时所说的话，与他们被欺骗时所说的话是不同的。因为当他们在教导他们的异端时，他们说的是他们被欺骗的事；但若他们说他们想他们不想的事，或他们不想他们想的事，他们说的是他们在说谎的事。在这件事上，任何人相信他们，即使没有发现他们，他自己也不会丧亡。因为如果当一个异端者谎称自己相信天主教的教义时，有人相信他是天主教徒，这并没有偏离天主教的准则：因为这对他是无害的；因为他在一个人的思想上犯了错（这个人隐藏时他无法判断），而不是在应当保存在他内的天主的信仰上犯了错。此外，当他们在教导他们的异端时，任何人若相信他们，认为那是真理，他将既分享他们的错误，也分享他们的惩罚。因此，当他们在编造他们邪恶的教义（他们被致命的错误所欺骗）时，那时相信他们的人就丧亡了；而当我们宣讲天主教的教义（我们持守正确的信仰）时，如果那人相信，那么失去的就找回了。但是，当他们是\OriginalTerm{普里西利安派}{Priscillianists}，为了不暴露他们的毒液，而谎称是我们中的一员时；我们中任何人若相信他们，即使他们没有被发觉，他自己仍坚持作一个天主教徒；另一方面，我们，如果为了发现他们而假称自己是普里西利安派，因为我们赞美他们的教义作为我们自己的，任何人若相信这个，要么会在这群人中被巩固，要么会立刻被转移到他们中间；但随后的一刻会带来什么？他们是否会被我们后来讲真话时从那里解放出来，因为当时他们被我们讲假话所欺骗？他们会愿意听一个他们亲身经历过说谎的人教导吗？谁能确切知道？谁能不知道这是不确定的？由此可以得知，天主教徒说谎以捕捉异端者，比异端者说谎以免被天主教徒发现，更为有害，或者温和地说，更为危险。因为，相信天主教徒说谎来试探人的，要么成为异端者，要么被巩固为异端者；但相信异端者说谎来隐藏自己的，并不停止作一个天主教徒。但为了使这一点更清楚，让我们举例来说明，并优先从您寄给我阅读的那些著作中选取例子。\par

5. 那么，让我们在眼前摆出一个狡猾的探子，他接近一个已认出是普里西利安派的人；他从主教狄克提纽开始，说谎地称赞他的生平（如果认识他），或他的名声（如果不认识他）；到此为止还算可以容忍，因为狄克提纽被认为曾是一位公教徒，并已从那个错误中纠正过来。然后，转向普里西利安（因为这在说谎的艺术中是下一步），他应当以尊敬的口吻提及那个不敬神且可憎的人，因他的邪恶罪行和罪过而被定罪！在这种尊敬的口吻中，如果这个为此类网罗所设的人碰巧不是一个坚定的普里西利安派，那么通过这番吹捧，他就会被巩固。但当探子继续谈论其他事情，说他怜悯那些人，因为黑暗的作者在如此黑暗的错误中召唤他们，以致他们不认识自己灵魂的荣耀和神圣始祖的光辉时：然后论及狄克提纽的书，那书被称为\WorkTitle{磅}，因为它从头到尾讨论了十二个问题，如同构成一磅的盎司，他将用这样的赞美来颂扬它，以致宣称这样一个\WorkTitle{磅}（其中包含着可怕的亵渎）比成千上万磅黄金更宝贵；的确，说谎者的这种狡猾扼杀了相信它之人的灵魂，或者，那灵魂已死，就在同样的死亡中沉没并压制它。但你会说：\Quote{之后它将被释放。}如果事情没有发生，要么因为介入某事阻止已经开始的事完成，要么因为异端思想的固执，虽然已经开始承认某些事，却又否认同样的事，那怎么办？尤其是，如果他发现受到一个陌生人的干涉，他将更加大胆地致力于用谎言隐藏自己的意见，因为他会更确定地知道，这样做是无可指责的，甚至以干涉者的例子为证。这，在一个认为用谎言隐藏真理是正确的人身上，我们有什么脸面指责，并敢于谴责我们所教导的事呢？\par

6. 因此，剩下的就是，普里西利安派根据他们异端邪恶的虚假，对天主、灵魂、身体以及其他所想的，我们毫不迟疑地以真实的怜悯加以谴责；但他们对用谎言隐藏真理的权利所想的，对我们和他们来说（愿天主禁止！）将成为一个共同教义。这是一个如此巨大的邪恶，即使我们这种通过谎言来捕捉并改变他们的企图如此成功，以致我们真的捕捉并改变了他们，也没有任何收益能补偿为了纠正他们而使我们与他们一起在错误上受损的损失。因为通过这个谎言，我们在这方面也变得乖戾，而他们只是被纠正了一半；因为他们认为为了真理而说谎是正确的，这是一个我们不予纠正的过错，因为我们学会了并教导同样的事，并且规定这样做是适宜的，以便我们能够达到对他们的修正。然而，我们并没有修正他们，因为他们的过错——他们认为用谎言隐藏真理是正确的——我们并未除去，反而当我们通过这样的过错寻找他们时，我们自己也变得有过错；我们也不知道如何相信他们（当他们悔改时），因为我们曾向他们（当他们执迷不悟时）说了谎；免得他们被捕获时，对他们所做的事，在捕获后也对我们做；不仅因为这是他们的习惯，而且因为我们这些被他们接近的人，也发现了同样的事。\par

7. 而且，更悲惨的是，即使在他们已经仿佛属于我们之后，他们也无法相信我们。因为如果他们怀疑，甚至在大公教义本身中，我们也在说谎，以隐藏我不知道的其他我们认为正确的事；当然，对于一个同样怀疑的人，你会说：\Quote{我只是在那个时候为了捕获你才这样做}；但当他问：\Quote{我怎么能知道你现在不是为了避免被我捕获而说谎呢？}时，你将如何回答？或者，难道有人能相信一个为了捕获而说谎的人，不会为了避免被捕获而说谎吗？你看到这个邪恶导致什么吗？那就是，不仅我们对他们，他们对我们，而且每个弟兄对每个弟兄都理所当然地变得可疑？于是，通过谎言所欲达到的目标（即信德可以被教导）所导致的，反而是在任何人身上都没有信任。因为如果我们说谎时甚至反对天主，那么人们在任何谎言中能发现什么如此巨大的邪恶，以致我们应当以各种方式躲避它，好像它是极其悲惨的事？\par

8. 但是现在请观察，与我们相比，普里西利安派的谎言尚且更可容忍，因为他们知道自己说话是欺诈的；我们却认为凭自己的谎言去拯救他们脱离那些因错误而腐朽的虚假事物是正当的。普里西利安派说：灵魂是天主的一部分，与祂具有同一本性和实体。这是一种重大而可憎的亵渎。因为由此推出，天主的本性可以被俘虏、被欺骗、被蒙蔽、被搅乱、被玷污、被定罪、受折磨。但如果那个为了从如此大恶中用谎言拯救一个人而这样说的人，也说出这话，那么让我们看看两个亵渎者之间有何区别。\Quote{很大区别}，你说：\Quote{因为普里西利安派这样说，也如此相信；但公教徒却不如此相信，虽然如此说。}那么，一个人是在无知中亵渎，另一个人是在有知中亵渎：一个违背知识，一个违背良心：一个具有思想虚假事物的盲目，但至少在那些事物上有意愿说出真实；另一个在暗中看见真理，却情愿说出虚假。\Quote{但一个}，你会说，\Quote{教导这个，为使人们成为他的错误和疯狂的参与者：后者说出它，为从那个错误和疯狂中拯救人们。}如今我在上面已经指出了这被人们认为将有益的事本身是多么有害：但同时，如果我们在这两者中衡量眼前的恶（因为公教徒从纠正异端者所寻求的未来之善是不确定的），谁犯罪更重？谁是在无知中欺骗人，还是谁是在有知中亵渎天主？无疑，哪一个是更严重的，那个以忧惧的虔敬将天主置于人之上的人就明白。再加上这一点：如果为了引导人们赞美天主而可以亵渎天主，那么无疑我们以自己的榜样和教训邀请人们不仅赞美，而且亵渎天主：因为我们谋图通过亵渎天主的行径将人们带到对天主的赞美中，如果我们果真把他们带来了，他们确实会学会不仅赞美，而且亵渎。这就是我们给予那些人的利益，我们通过并非无知而是有知的亵渎将他们从异端中拯救出来！而宗徒将人交给撒殚自己，为使他们学习不再亵渎；我们却努力将人从撒殚手中救出，为使他们学习不是无知地而是有知地亵渎。而且我们这些他们的师傅，给自己带来如此大的祸害，以致为了捕捉异端者，我们首先成为——这是确定的——亵渎天主的人，以便为了拯救他们——这是不确定的——能够成为祂真理的教师。\par

9. 因此，当我们教导我们的人亵渎天主，以便普里西利安派相信他们是自己人时，让我们看看他们自己为了让我们相信他们是自己人而说谎时说了什么恶事。他们按我们的心意诅咒普里西利安，憎恶他；他们说灵魂是天主的受造物，不是一部分；他们斥责普里西利安派的虚假殉道；他们极力赞扬那些揭露、攻击、倾覆该异端的公教主教，等等。看，当他们说谎时，他们自己说的是真话：并不是说谎言本身可以同时是真实的；而是当他们在某件事上说谎时，在另一件事上说了真话：因为当他们在说属于我们这件事上说谎时，对公教信仰说了真话。因此，他们为了不被发现是普里西利安派，以说谎的方式说出真理：但我们为了发现他们，不仅说谎让他们相信我们属于他们；而且我们还说出虚假的事，即我们知道属于他们错误的事。因此，当他们希望被认为属于我们时，他们所说的话一部分是假的，一部分是真的；因为他们属于我们是假的，但灵魂不是天主的一部分是真的：但当我们希望被认为属于他们时，我们所说的话两种都是假的，既假称我们是普里西利安派，又假称灵魂是天主的一部分。那么，当他们隐藏自己时，他们赞美天主，并不亵渎；当他们不这样做，而是说出自己的情感时，他们不知道自己在亵渎。所以，如果他们皈依公教信仰，他们可以安慰自己，因为他们能说出宗徒所说的话：他在其他事中曾说：\BibleQuote{我先前是个亵渎者；但}，他说，\BibleQuote{我蒙了怜悯，因为我是无知而做的。}相反地，我们为了使他们向我们敞开，如果我们将这话说出来，好像是为欺骗和捕捉他们而说的一个正当谎言，那么我们就确实既声称我们属于亵渎的普里西利安派，又为使他们相信我们，而在没有无知借口的情况下亵渎。因为一个公教徒，他想通过亵渎被当作异端者，就不能说：\Quote{我是无知而做的。}\par

10. 我兄弟，在这样的情形中，常常应当怀着恐惧记起：\BibleQuote{凡在人前否认我的，我必在我天父前否认他。}或者，真的在人前否认基督不是否认基督吗？在普里西利安派面前否认祂，为的是当他们隐藏自己时，一个人可以通过一个亵渎的谎言揭露并捕捉他们？但是，请问，有谁怀疑基督被否认了，当我们不按祂的真实所说祂不是，而按普里西利安派所相信的所说祂是？\par

11. \Quote{但是，隐藏的狼}，你会说，\Quote{披着羊皮，暗中严重危害主的羊群，我们无法以其他方式发现他们。}那么，在编造这种用谎言追捕他们的方法之前，\OriginalTerm{普里西利安派}{Priscillianists}是如何为人所知的？他们那位无疑更加狡猾因而也更加隐蔽的创始人，又是如何在床上被找到的？如此众多、如此重要的人物是如何被揭露和定罪，而其余无数的人，部分得到纠正，部分似乎得到纠正，并在教会的怜悯中被聚集到她的羊圈里？因为主在怜悯时赐予许多途径，使我们能够发现他们；其中两条途径比其他途径更幸福：即要么是他们曾想引诱的人，要么是他们已经引诱了的人，当他们悔改并皈依时，会指认他们。这一点更容易实现，如果他们的邪恶错误不是通过欺骗的伎俩，而是通过真实的推理被推翻的话。在撰写这些文字时，你应当付出辛劳，因为天主赐予你这样的能力：这些有益的著作摧毁了他们疯狂的乖谬，这些著作越来越为人所知，并由公教徒，无论是在集会中演讲的主教，还是任何充满天主热诚的勤勉之人，到处传播，这些将成为神圣的网罗，在其中他们可以被真实地捕获，而不是通过谎言被追捕。因为这样被捕获，要么他们自愿承认自己曾经是什么，并认识其他属于邪恶团体的人，他们要么仁慈地纠正，要么慈悲地告发。或者，如果他们羞于承认自己通过长期伪装所隐藏的东西，那么天主隐藏的手会医治他们，使他们痊愈。\par

12. \Quote{但是}，你会说，\Quote{如果我们假装与他们一样，就能更容易地看穿他们的伪装。}如果这是合法或有益的，基督本可以指示他的羊群，让他们披着狼皮来到狼群中，并通过这种欺骗手段发现他们；但他没有这样说，甚至在他预言他将把羊群送入狼群中时也没有说。但你会说：\Quote{那时他们不需要被调查，因为他们是最明显的狼；但他们的撕咬和凶残必须忍受。}那么，当预言后来的时期，他说凶残的狼会披着羊皮而来时，难道没有空间给出这个建议并说：你们也要为了发现他们而披上狼皮，但里面仍然是羊吗？他并没有这样说；而是当他说了\BibleQuote{许多人会披着羊皮来到你们这里，里面却是凶残的狼}之后，他接着说，不是\Quote{通过你们的谎言}，而是\BibleQuote{通过他们的果实你们会认出他们}。我们必须通过真理来提防谎言，通过真理来捕获，通过真理来消灭谎言。但愿我们绝不通过故意亵渎来战胜无知者的亵渎；但愿我们绝不通过模仿来防备欺骗者的邪恶。因为如果我们为了防备他们而拥有他们，我们怎能防备他们呢？因为如果为了抓住一个无意中亵渎的人，我故意亵渎，那么我所做的比我所抓的更坏。如果为了找出一个无意中否认基督的人，我故意否认他，那么我这样找到的人将会跟随我走向毁灭，因为我为了找出他而先毁灭了自己。\par

13. 或许情况是这样：以这种方式谋划找出\OriginalTerm{普里西利安派}{Priscillianists}的人，并没有否认基督，因为他的口说出他心中不相信的话？好像真的（我上面也说过）当说\BibleQuote{人心里相信，就可成义}时，又白白地加上了\BibleQuote{口里宣认，就可获得救恩}吗？难道不是几乎所有的在迫害者面前否认基督的人，心中都持有他们对基督的信仰吗？然而，由于他们没有用口宣认以至于救恩，他们就灭亡了，除非他们通过补赎重新活过来。谁能如此荒谬，认为宗徒伯多禄在否认基督时，他嘴上说的就是他心中想的呢？当然，在那次否认中，他内心持守真理，而口出谎言。那么，为什么他用眼泪洗去他用口说出的否认呢？如果心里相信就足以得救，那又何必呢？他为何在心中说真话，却用痛苦哭泣惩罚他口中说出的谎言？除非因为他看到这是一个巨大的、致命的恶：当他心里相信以至于成义时，他口里却没有宣认以至于救恩。\par

14. 因此，所记载的\BibleQuote{那在自己心里说真话的人}，不应被理解为：真理存于心中，口里却可以说谎。之所以这样说，是因为一个人有可能用口说出对他毫无益处的真理，如果他没有把它放在心里，也就是说，如果他自己不相信他所说的话；就像异端者，尤其是这些\OriginalTerm{普里西利安派}{Priscillianists}所做的那样，他们确实不信公教信仰，却还是说出它，使他们被相信是我们的人。因此，他们口里说真理，心里却不。正因如此，他们要与那所记载的\BibleQuote{他在心里说真话的人}区分开来。如今，公教徒在心里说这真理，因为他如此相信，所以也应当用口说，以便宣讲它；但反对它，无论在心上还是口上都没有虚假，以便他一方面心里相信以至于成义，另一方面用口宣认以至于得救。因为在那篇圣咏中，在说了\BibleQuote{那在自己心里说真话的人}之后，紧接着又加上了\BibleQuote{那没有用舌头行诡诈的人}。\par

15. 至于宗徒所说：\Quote{\BibleQuote{弃绝谎言，各人与近人说实话，因为我们彼此是肢体}}，我们万不可如此理解，仿佛他允许我们与那些尚未成为基督身体肢体的外人说话时可以说谎。这话之所以如此说，是因为我们每人都应将他人视为我们所希望他成为的那样，即使他尚未成为那样；正如主将怜悯那受伤者的外邦撒玛黎雅人视为他的近人。因此，与我们打交道的人应被视为近人，而非外人，因为我们关心的是他不至于成为外人。即便因他尚未分享我们的信德和圣事，有些真理事必须向他隐瞒，但这绝不是应当向他撒谎的理由。\par

16. 因为在宗徒时代，有些人并非以诚实的心传播真理，即是说，他们心怀诡诈。关于这些人，宗徒说他们传播基督不是出于纯洁，而是出于嫉妒和纷争。为此，当时有些人虽不怀纯洁之心传播真理，仍被容忍；但从未有人因以纯洁之心传播谎言而受称赞。最后，宗徒论到他们说：\Quote{\BibleQuote{或是假意，或是诚心，无论如何，基督究竟被传开了}}。但他绝不会说：为使基督以后被传扬，应先否认祂。\par

17. 因此，虽然确实有许多方法可以在不诋毁天主教信仰、不赞扬异端邪说的前提下，搜寻隐藏的异端者；但若除此之外别无他法将异端邪说从其洞穴中引出，而必须使天主教口舌偏离真理的正道，那么，宁可让异端隐藏，也不可让天主教口舌失足；宁可让狐狸躲在洞中不被看见，也不可为了捕捉它们而使猎人坠入亵渎的深坑；宁可让普里西利安派的背信行为被真理的帷幕所掩盖，也不可让天主教信仰因担心受说谎的普里西利安派赞扬而被有信德的天主教徒否认。因为如果谎言（且不论何种谎言），尤其是亵渎的谎言，因其出于揭露隐藏异端者的意图而被视为正义，那么照此推理，若以同样意图交换，则也可能产生贞洁的奸淫。不妨设想：在一群淫荡的普里西利安派中，有一个女人看中了一位天主教的若瑟，她答应若瑟，若他同她同寝，她便出卖他们的秘密巢穴，且若能确定若瑟应允她，她定会履行诺言。我们是否该认为这事应当去做？还是我们当明白，绝不可为此种货色付出如此代价？那么，为何我们为了捕捉异端者，不通过肉体犯奸淫的放荡来驱赶他们，却以为通过口舌犯亵渎的淫乱来驱赶他们是正当的呢？因为要么可以同样理由为两者辩护，说这些事之所以并非不义，是因为做这些事出于发现不义之人的意图；要么健全的道理既不许我们为了发现异端者而与不洁的女人发生关系（即使仅限于身体，而非内心），那么同样，也不许我们为了发现异端者而通过言辞（即使仅限于声音，而非内心）去宣扬不洁的异端或亵渎纯洁的天主教会。因为那应掌管人一切下等行动的理智之主权，若做了不应做的事（无论通过肢体还是言辞），就免不了应受的责难。尽管即使通过言辞行事，也是通过肢体行事，因为舌头是造言的肢体；而且我们任何肢体行为唯有先在心中孕育才能产生，甚或当我们内心思想并同意时已产生，然后通过一个肢体在外在行为中生出。因此，若有人说某事不是出于内心意愿，但若没有内心决定去做，它就不会发生，那么内心就不能以此推卸责任。\par

18. 所作之事的因由、目的和意向的确有很大区别；但是那些明显是罪的事，不论出于何等善因、看似何等善果、自称何等善意，都不可去做。因为人的行为若非本身为罪，则或因善因而善，或因恶因而恶；例如，因怜悯之心且怀有正确信德而施食给穷人，是为善工；为生养子女而履行夫妻之礼，且心怀信德以生育可受洗的后代，也是善工。此类行为按其因由或善或恶，因为同样的行为若出于恶意，便成为罪：例如，为炫耀而施食穷人，或为纵欲而行夫妻之礼，或生养子女并非为献给天主，而是献给魔鬼。但是，当行为本身为恶时，诸如偷窃、奸淫、亵渎等，有谁能说可因善因而行，以致它们要么不是罪，要么更荒谬地只是罪呢？有谁能说：\Quote{我们为了周济穷人，就去偷窃富人；或者我们作伪证，尤其是不伤害无辜，反而解救那些将被法官定罪的罪人？}因为作这等伪证可带来两件好事：一边得钱养活穷人，一边欺骗法官使人不受罚。甚至在遗嘱的事上，如果我们能，为何不抑制真遗嘱而伪造假遗嘱，使遗产或遗赠不落入那些不行善事的不配之人手中，反而落入那些用以喂饱饥饿者、衣着赤身者、款待旅客、赎回俘虏、建造教堂的人手中？既然为了这些善事，那些恶事根本就不是恶，为何不为这些善事而行那些恶事？更有甚者，倘若淫荡富有的妇人还可能为她们的情夫带来财富，那么一个心怀慈悲的人为何不可以从事这些行当，以便有财源周济穷人，而不听宗徒的话：\BibleQuote{偷窃的，不要再偷；反而更要劳苦，亲手做正当的事，好能周济有急需的人。} \BibleRef{弗 4:28} 如果不仅偷窃，而且伪证、奸淫和一切恶行，若为了行善而做，就不是恶而是善，那有谁能说这样的话呢？除非是想要颠覆人间一切风俗和法律的人。因为有什么最凶恶的罪行、最污秽的丑事、最不敬的亵渎，不能说是做得正确和正义的？不仅不受罚，甚至可荣耀地施行，不仅不怕惩罚，还可能盼望赏报？只要我们一旦承认，在人的一切恶行中，不应问所作的是什么，而应问为何而作；并且，那些为善因而作的事，甚至不应被判定为恶。但是，如果正义理应处罚一个贼，尽管他自称并证明他为富人之多余而取之，以给穷人之所需；如果正义理应处罚一个伪造者，尽管他证明他篡改别人的遗嘱，是为了让那将广施哀矜的人继承，而不是那不做哀矜的人；如果正义理应处罚一个奸夫，尽管他证明他是出于怜悯而犯奸，为了通过与他行淫的女子解救一个人免于死亡；最后，更进一步接近我们的问题，如果正义理应处罚一个意图与知情于\OriginalTerm{普里西利安派信徒}{Priscillianists}丑行的女子行淫，以便潜入他们秘密中的人；那么，请问，宗徒说：\BibleQuote{不要将你们的肢体交给罪恶，作不义的武器；} \BibleRef{罗 6:13} 那么，我们的舌头、我们的整个口、我们的发音器官有何亏负我们，以至于我们要将它们交给罪恶作武器，而且犯下如此大罪，为了擒获并解救那些无知亵渎的普里西利安派信徒，我们自己却毫无无知之借口而亵渎我们的天主？\par

19. 有人会说：\Quote{那么，任何贼都与那位出于怜悯而偷窃的贼同等吗？}谁会这样说？但是，在这两者之间，不能因为一个更坏就认为另一个是好的。出于贪心而偷窃的贼比出于怜悯而偷窃的贼更坏；但若所有偷窃都是罪，我们就必须远离一切偷窃。因为谁能说人可以犯罪，即使一个罪应受罚，另一个罪只属小罪？但目前我们问的是：若人做这事或那事，谁不会犯罪或谁会犯罪？而不是谁犯罪更重或更轻。因为即使是偷窃，法律对它的处罚也比对淫秽之罪的处罚轻；然而两者都是罪，虽然一个较轻，一个较重；以至于出于贪欲而犯的偷窃被认为比出于行善而犯的淫秽之罪更轻。也就是说，在同类罪中，那些似乎出于善意而犯的罪比同类其他罪更轻；然而，当与另一类罪比较时，就类别本身而言更轻的罪，却可能显得更重。出于贪婪而偷窃比出于怜悯而偷窃更重；同样，出于淫欲而行淫比出于怜悯而行淫更重；然而，出于怜悯而犯奸淫可能比出于贪婪而偷窃更重。我们现在关心的不是孰轻孰重，而是什么是罪，什么不是罪。因为没有人可以说：在明显是罪的事上，犯罪是义务；但我们说：若罪是如此这般犯下的，是义务要宽恕或不宽恕。\par

20. 但是，必须承认，某些补偿性的罪过，在人心中造成如此大的困惑，以至于它们甚至被认为值得赞美，而更应该被称为正确行为。因为，如果一位父亲将自己的女儿交给不敬神的人行淫，谁能怀疑这是大罪呢？然而，曾出现过一个案例，一位义人认为他有责任这样做，当索多玛人以邪恶的色欲猛烈冲向他的客人时。因为他说：\BibleQuote{我有两个女儿，还未认识男人；我把她们领出来，你们可以任意待她们；只是对这些人不可做坏事，因为他们已进入我屋顶的遮盖之下。}我们对此该说什么？难道我们不如此憎恶索多玛人企图对义人的客人所做的邪恶，以至于无论做什么，只要这事不发生，他就应认为是正确的吗？行此事者的身份也极大地打动了我们，他因义德之功劳而得以从\Place{索多玛}得救，以至于说，既然女子遭受淫乱比男子更轻，那么甚至属于那位义人的义德，他为自己的女儿们选择了此事而非为他的客人；不仅在心里愿意，而且还以言语提出，如果他们同意，就准备在行动上实现。但是，如果我们为罪恶开辟了这条道路，即我们要犯较小的罪，以免他人犯更大的罪；那么，通过一个宽广的边界，不，更确切地说，根本没有边界，而是撕裂并移除所有界限，在无限的空间里，所有的罪都将进入并作王。因为，当它被定义，一个人要犯较小的罪，以免另一个人犯更大的罪；那么，当然，通过我们犯偷盗，可以防范他人的淫乱，以及通过淫乱防范乱伦；如果任何不敬神的行为甚至比乱伦更坏，那么乱伦也将被宣称为我们应该做的，如果这样可以阻止他人不犯那不敬神的行为的话：在每一种罪行中，无论是偷盗换偷盗，淫乱换淫乱，乱伦换乱伦，都将被认为是应该做的：我们自己的罪为别人的罪，不仅是以小换大，甚至如果达到最高最坏的，以少换多；如果事情的紧迫如此转变，否则他人不会停止犯罪，除非通过我们犯罪，尽管是小一些的犯罪，但仍然是犯罪；以至于在任何情况下，当一个拥有这种能力的敌人说：\Quote{除非你作恶，我将更作恶，或者除非你行这恶，我将行更多这样的恶}，那么我们必须看似在自己身上接受邪恶，如果我们希望阻止（他人）行恶的话。这样明智，岂不是失去理智，或者更确切地说，是彻底疯狂？我自己的罪孽，而非他人的，无论是对我还是对他人所犯的，才是我必须提防被定罪的事。因为\BibleQuote{犯罪的灵魂，它必死。}\par

21. 那么，如果为了不让别人犯更重的罪——无论是针对我们还是针对任何人——而去犯罪，无疑我们不应该；我们需要考虑\Person{罗特}所做的，它是一个我们应当效法的榜样，还是一个我们应当避免的榜样。因为，似乎更值得深入观察和注意，当一场如此可怕的邪恶，来自索多玛人最凶残的不敬，即将降临到他的客人身上时，他想要避开却无法做到，以至于那位义人的心被扰乱到如此程度，以至于他愿意做那样的事——不是带着迷蒙情绪的人的恐惧，而是天主的法律在其平静安宁中（如果我们咨询它），会大声宣告：不可做！并会命令我们，要如此谨慎以防自己犯罪，以至于不因惧怕他人任何罪恶而犯罪。因为那位义人，因惧怕他人的罪恶（这只能玷污那些同意的人），如此心烦意乱，以至于没有注意到自己的罪，即他愿意将女儿们交给不敬神之人的情欲。这些事，当我们在圣经中读到的时候，我们不能因为相信它们发生过，就因此相信它们应当被做；免得我们在不加分辨地遵循先例时违反诫命。或者，真的，因为\Person{达味}发誓要杀死\Person{纳巴耳}，后来出于更慎重的仁慈而没有实施，我们因此就说他应当被效法，以便我们可以发誓要做一件事，之后又看到它不应该做吗？但是，正如恐惧扰乱了前者，使他愿意牺牲自己的女儿，同样，愤怒也扰乱了后者，使他鲁莽地发誓。总之，如果我们被允许询问他们两人，请他们告诉我们为什么做这些事，一人可能会回答：\Quote{恐惧和战栗临到我身，黑暗笼罩了我；}另一人也可能会说：\Quote{我的眼睛因愤怒而困扰；}因此我们不应奇怪，无论是在恐惧黑暗中的前者，还是眼睛困扰的后者，都没有看到当看的事，以免他们做出不当做的事。\par

22. 对\Person{圣达味}，确实更可公正地说，他不该发怒；即使对那忘恩负义、以恶报善的人也是如此；但若他作为人，怒气悄然袭来，他不该让怒气如此得胜，以致他发誓要做一件事，要么因屈服于愤怒而做，要么因违背誓言而放弃不做。但对另一位，置身于\Place{索多玛}人情欲狂乱之中，谁敢对他说：\Quote{你的客人，在你自己的家中，你用强暴的好客迫使他们进入，如今被淫秽的男人抓住，被强行玷污并像女人一样被肉体认识，你却不害怕丝毫，毫不关心，毫不畏惧、不惊恐、不战栗？} 什么人，即使是那些恶徒的同伴，敢对这位虔诚的主人说这样的话？但无疑，最正确的话是：\Quote{尽你所能，使那件你理当惧怕的事不发生：但不要让你这惧怕驱使你去做一件这样的事：如果你的女儿们愿意这事发生在她们身上，她们将因你而与索多玛人行邪恶；如果不愿意，她们将因你而遭受索多玛人的强暴。不要犯下你自己的大罪，同时却惧怕别人的更大罪行；因为无论你自己的罪与别人的罪之间的差别有多大，这罪将是你的，那是别人的。} 除非有人为这人辩护时如此强词夺理，说：\Quote{既然受损害比施损害更好，而那些客人不是要施损害，而是要受损害，那位义人便选择让他的女儿们受损害，而不是他的客人，因为他是女儿的主人，行使了权力；他知道，如果事情发生，在她们身上也不会有罪，因为她们只是承受那些犯罪的人，并不认同他们，所以她们自己没有罪。总之，她们并没有献出自己（尽管是比男性更好的女性）去代替那些客人被肉体认识，以免自己成为有罪，不是因承受别人的情欲，而是因自己意志的同意；父亲也没有允许这事发生在他自己身上，当他们试图这样做时，因为他不想把客人交给他们（尽管若只对一人而非两人做，邪恶会小一些）；但他尽可能抵抗，以免自己也因任何自己的赞同而被玷污，即使别人的情欲狂乱以体力得胜，只要他不赞同，也不会玷污他。既然女儿们没有犯罪，他也没有在她们身上犯罪，因为他并没有使她们犯罪，如果她们违背意志被强暴，只是承受犯罪者。就好像他应该献出他的奴仆让暴徒鞭打，以免客人在他的家里遭受鞭打之害。} 我不会辩论此事，因为辩论起来很长：是否主人可以正当地运用他对奴仆的权力，使一个无辜的奴仆被鞭打，以免他的无辜朋友在家里被凶恶的坏人鞭打？但肯定，关于\Person{达味}，说他应该发誓去做一件他后来会明白不该做的事，这是绝对不正确的。可见，我们不该把我们读到圣洁或义人所做的一切都当作道德规范，而应由此学习\Person{宗徒}的那句话多么广泛，甚至适用于哪些人：\BibleQuote{弟兄们，若有人一时不慎犯过，你们属灵的人就该以柔和的心神挽回这样的人；同时也要小心，免得自己也受诱惑。} 一时不慎犯过，要么是当时没有看见什么是正确的事而做，要么是看见了却被打败；也就是说，一个罪被犯下，要么因为真理被隐藏，要么因为软弱所迫。\par

23. 但在我们的一切行为中，即使是善人，在\OriginalTerm{补偿之罪}{compensative sins}的事情上也大感困扰；以至于这些罪不被认为是罪，如果它们有那为此而行的原因，并且在这些原因中，若不做，反而更像是罪。尤其在关于谎言的事上，在人的意见中已到了如此地步：那些谎言不被算为罪，反而被认为是正当的，当一个人为那个被欺骗对他有利的人的利益而说谎，或为了避免某人伤害别人，而那人若不靠谎言打发走似乎可能伤害。为辩护这类谎言，人们认为有极多圣经的例子支持。然而，隐藏真理与说出谎言并不是一回事。因为每个说谎的人都希望隐藏真实的事，但不是每个希望隐藏真实的事的人都说谎。因为通常我们不是以说谎而是以沉默来隐藏真理。因为主没有说谎，当祂说：\BibleQuote{我还有许多事要告诉你们，但你们现在不能承担。} 祂对真实的事保持沉默，而不是说虚假的事；因为祂判断他们不适于听那些真理。但如果祂没有向他们指明这一点，即他们不能承受祂不愿说的话，那么祂确实隐藏了一些真理，但这样做是否正当，我们或许不知道，或者没有如此伟大的榜样来坚定我们。因此，那些主张有时可以说谎的人，并不恰当地提到\Person{亚巴郎}关于\Person{撒辣}的事，他说她是他的妹妹。因为他没有说：她不是我的妻子，而是说：\BibleQuote{她是我的妹妹；}因为她实际上血缘如此之近，以至于可以不说谎地称为妹妹。后来，在把她交还给那个带走她的人之后，他也确认了这一点，回答他说：\BibleQuote{诚然她是我的妹妹，是父亲的，不是母亲的；}即父系亲属，不是母系。因此，他隐瞒了一些真实的事，没有说出任何虚假的事，当他未提及妻子而提到妹妹时。他的儿子\Person{依撒格}也做了同样的事：我们知道他也娶了一位近亲。因此，当以沉默保留真实的事时，不是谎言；但当以言语说出虚假的事时，才是谎言。\par

24. 关于\Person{雅各伯}，他在母亲指使下所做的，看似欺骗了父亲，但若以勤奋和信德来审视，那并非谎言，而是一个奥迹。我们若将这些称为谎言，那么一切比喻和旨在象征任何事物的形象，凡是不可按其字面意义理解，而须由此及彼领悟的，也都应称为谎言了；这断然不可。因为如此思想的人，对于如此众多的比喻表达，也可能将这一诽谤加诸其上，以至于连所谓的隐喻（即一字从其本意不当转用于另一对象）也可据此称为谎言。因为当他说到波浪般的麦田、结出宝石的葡萄藤、青春的绽放、雪白的头发时，毫无疑问，波浪、宝石、绽放、白雪，由于我们在那些对象中找不到这些词语从其本处转用来的东西，便会被这些人当作谎言。而基督是磐石，犹太人的石心；基督是狮子，魔鬼也是狮子，以及无数类似的说法，都将被说成是谎言。不仅如此，这种比喻表达甚至触及所谓\OriginalTerm{反义法}{antiphrasis}，例如说某物丰盛实则不存在，说某物甜美实则是酸涩；\Quote{\OriginalTerm{树林因其不发光，命运女神因其不怜悯}{lucus quod non luceat, Parcae quod non parcant}}。属于这类的，有圣经中魔鬼对主论及圣约伯时说的话：\BibleQuote{他岂不当面祝福你？}其意思是\Quote{诅咒}。纳波特被诬告的罪名也用了这个词；据说他\BibleQuote{祝福了国王}，也就是诅咒了国王。所有这些说话方式，若将形象化的言语或行动视为说谎，就都会被算作谎言。然而，若借他物以表此物的说法，指向对真理的理解，便不是谎言；那么，不仅\Person{雅各伯}为得祝福对父亲所做所说的，而且\Person{若瑟}仿佛嘲弄他的兄弟所说的话，以及\Person{达味}装疯，都必不被判为谎言，而是预言性的言论和行动，指向对真理之事的理解；它们仿佛披着比喻的外衣，目的是操练虔敬探求者的心智，免得因暴露在表面而变得廉价。纵然我们从别处——那里说得公开明确——所学到的事物，当它们从隐藏之处被取出时，因我们（在某种意义上）发现了它们，便得以更新，且因更新而甘甜。它们并非因以这些方式隐藏而吝啬于学习者；而是以一种更吸引人的方式呈现，仿佛被收回，使它们被更热切地渴望，渴望之后被更欣喜地找到。然而所说的是真理，而非虚假；因为无论言辞或行为所象征的都是真理，而非虚假；所象征的正是所说的。它们之所以被当作谎言，只因人们不理解所象征的真理正是所说的，而认为所说的是虚假的。为使这点更明白，请留意\Person{雅各伯}所做的这件事。无疑，他用山羊羔皮遮盖肢体；若寻求直接原因，我们会认为他撒谎了；因为他这样做，是为了让人以为自己不是本人：但若此举指向它真正象征的对象，则山羊羔皮象征罪恶；用皮遮盖自己的那位，是承担了并非自己的、而是他人的罪恶的那一位。因此，真实的象征决不能正当地称为谎言。正如在行为上，在言语上也是如此。即当他父亲对他说：\BibleQuote{你是谁，我的儿子？}他回答：\BibleQuote{我是厄撒乌，你的长子。}这若指向那对孪生兄弟，将看似谎言；但若指向这些言行所象征的那一位，则在此应理解为那在祂的身体——教会——内的那一位，祂论及此事时说：\BibleQuote{你们将看见亚巴郎、依撒格、雅各伯及众先知在天国里，而你们被赶出去。他们将从东方、西方、北方、南方而来，并坐席于天国；看，有最后的将成为最先的，有最先的将成为最后的。}既然如此真实、如此真实地被象征，这里有什么应被视为是虚假地做或说的呢？因为当所象征的并非虚假之事，而是真实之事——无论是过去、现在还是将来——毫无疑问，那是真实的象征，而非谎言。但在这种预言象征方面，要剥去外壳探寻所有细节则过于冗长；在当中真理取得了胜利，因为正如它们借象征被预示，后来也因应验而变得清晰。\par

25. 如今在这篇论述中，我并未承担更多属于你的工作，因为你已揭露了普里西利安派\OriginalTerm{普里西利安派}{Priscillianists}的藏身之处，即关乎他们虚假而乖谬的教义；以免人们以为这些教义被调查的方式，仿佛它们值得被教导，而非被驳斥。因此，请你更加努力去击垮并推翻它们，正如你曾努力去揭露并曝光它们一样；免得当我们希望揭露那些行虚假之人的时候，我们却允许虚假本身，仿佛不可战胜一般，固守阵地；我们本应甚至在这些隐藏的异端者的心中摧毁虚假，而不是因宽容虚假而找出那行虚假的欺骗者。此外，在他们那些应当被颠覆的教义中，有一条他们所教条化的，即：为了隐藏宗教，虔诚之人应当说谎，甚至达到这样的程度，以至于不仅关于其他不涉及宗教教义的事，而且关于宗教本身，也应当说谎，以免宗教暴露给外人；也就是说，人可以否认基督，以便在基督的敌人中间秘密地作基督徒。我恳求你，也要推翻这条邪恶而渎神的教义；他们为了支撑它，在论证中从圣经收集见证，以显得谎言不仅应被宽恕和容忍，甚至应被尊崇。因此，在驳斥那可憎的教派时，你有责任证明，那些圣经见证应当如此接受：要么你教导那些被认为是谎言的事，若按它们应当被理解的方式理解，便并非谎言；要么那些明显是谎言的事不值得效法；要么最终，至少关于那些涉及宗教教义的事，绝不可说谎。因为这样，他们才真正从根基上被推翻，即他们藏身之处被推翻：他们在这件事上最不适合我们效法，最应当被回避，因为他们为了隐藏自己的异端而自称是说谎者。这正是他们身上必须首先被攻击的，这仿佛是他们合宜的堡垒，必须以真理的打击来击碎并推翻。我们也不应给他们另一个他们所没有的藏身之处，让他们可以逃入其中，即，也许当他们被那些他们曾试图引诱却未能成功的人所背叛时，他们会说：\Quote{我们只是想试探他们，因为明智的公教徒曾教导，为了查出异端者，这样做是对的。}但是，有必要以更为恳切地请求你的恩惠来说明，为何在我看来，对于那些想用神圣圣经作为他们谎言的辩护者的人，这种三分法的辩论方式是合适的；即：首先，表明圣经中某些被认为是谎言的事情，若正确理解，并非如其所被认为的那样；其次，若那里有任何明显的谎言，它们不值得效法；第三，与所有那些认为有时说谎属于善人责任的人的意见相反，必须坚持在宗教教义上绝不可说谎。因为这就是我此前不久推荐并某种程度上嘱咐你去追求的这三件事。\par

26. 那么，为要表明圣经中那些被认为谎言的事，若正确理解，并非如所认为的那样，但愿你不要认为这不足以反驳他们，即他们并非从宗徒的著作，而是从先知书中找到谎言的先例。因为他们指名提及的那些人，其中每个人都说了谎，这些事都记载于那些书中，书中不仅记载了言语，也记载了许多具有比喻意义的行为，因为那些行为也是在比喻意义上完成的。但在比喻中，那些看似谎言的话，若好好理解，便会发现是真理。然而，宗徒们在他们的书信中是以另一种方式说话的，\WorkTitle{宗徒大事录}也是以另一种方式写成的，因为此时新约已经启示出来，它曾隐藏在先知的那些比喻中。简言之，在所有那些宗徒书信中，以及在那部以正典真理叙述他们事迹的大书中，我们找不到任何一个说谎的人，以至于这些人可以把他作为谎言的先例提出来。因为伯多禄和\Person{巴尔纳伯}的那种伪善，他们借此强迫外邦人犹太化，理应受到责备并改正，以免当时造成伤害，也不至于在后来成为可效法的事。因为当宗徒\Person{保禄}看见他们不按福音的真理正直行事时，就在众人面前对\Person{伯多禄}说：\BibleQuote{你既然是犹太人，却像外邦人而不像犹太人一样生活，怎么还能强迫外邦人犹太化呢？}但至于他自己所做的，为了借着遵守并实行某些按犹太习俗的律法规条，以表明他并非法律和先知的敌人，我们绝不可相信他是以说谎者的身份这样做的。事实上，关于此事，他的意见是众所周知的，据此确定：当时信了基督的犹太人不应被禁止遵守祖传的传统，外邦人成为基督徒后也不应被迫遵守这些；目的是使那些众所周知为由天主所命的神圣礼规，不被视为亵渎而回避；也不被视为如此必要，以至于新约启示后，若没有它们，任何归向天主的人便不能获得救恩。因为有一些人虽已接受了基督的福音，却仍这样想并如此宣讲；\Person{伯多禄}和\Person{巴尔纳伯}曾假意同意他们，因而强迫外邦人犹太化。因为宣讲它们如此必要，仿佛在接受福音之后，没有它们便不能在基督内获得救恩，这本身就是一种强迫。某些人的错误这样假设，\Person{伯多禄}的恐惧这样假装，\Person{保禄}的自由这样击破。因此，他所说：\BibleQuote{我成了众人中的一切，为要拯救众人。}他是借着与人同苦而这样做的，并非借着说谎。因为每一个人，当他以同样的怜悯施助时，他愿意自己若处于同样的困境也受到这样的施助，他就仿佛成了他所愿意救助的那个人。因此他仿佛成了那个人，并不是因为他欺骗他，而是因为他把自己当作他。由此有宗徒前面引用的话：\BibleQuote{弟兄们，如果有人一时不慎犯了过失，你们这些属灵的人，要用温和的精神去纠正他；但你自己要小心，免得也受诱惑。}因为如果因他所说：\BibleQuote{对犹太人，我就成为犹太人；对在法律下的人，我就成为在法律下的人。}便认为他以说谎的方式接受了旧约的礼仪，那么他也应该以说谎的方式接受外邦人的偶像崇拜，因为他说过对没有法律的人，他就成为没有法律的人；然而他绝对没有这样做过。因为他从未在任何地方向偶像献祭，或崇拜那些偶像，反而作为基督的殉道者自由地表明它们应受憎恶和远离。因此，这些人不能从宗徒的任何行为或言论中，举出可效法的谎言例子。那么，他们自以为能从先知的言行中找到依据，其唯一原因就是他们把预示性的比喻当作谎言，因为这些比喻有时与谎言相似。但当它们被关联到那些所象征的事物时，这些行为或言语就是为象征那些事物而如此行或说，它们就被发现是充满真理的象征，因此绝不是谎言。谎言，即有意欺骗的错误象征。但那不是错误的象征，尽管一件事被另一件事所象征，如果正确理解，所象征的事物是真实的。\par

27. 在福音中，甚至我们的救主也有这类事，因为先知的主也屈尊自己成为先知。例如，关于患血漏的妇人，祂说：\BibleQuote{谁摸了我？}；关于\Person{拉匝禄}，祂问：\BibleQuote{你们把他安放在哪里？}祂发问，仿佛不知道祂其实完全知道的事。祂因此假装不知道，为要借着祂看似无知来象征别的什么；既然这个象征是真实的，它肯定不是谎言。因为无论是患血漏的妇人，还是已死四天的人，都象征了某些人，连那知道万物的祂也在某种方式上不认识他们。因为那妇人预表外邦人的百姓，关于他们早有预言：\BibleQuote{一个我不认识的民族服事了我。}；而\Person{拉匝禄}从活人中被移去，仿佛在那个地方以象征性的相似躺卧，那里正是那一位所躺卧的地方，祂的声音是：\BibleQuote{我从你眼前被抛弃了。}正是为此，基督仿佛不知道她是谁，也不知道他被安放在哪里，借着祂询问的话，塑造了一个比喻，并借着真实的象征，完全排除了谎言。\par

28. 因此，你所提到的他们所述之说，即主耶稣复活后，与两位门徒同行；当他们临近他们要去的村庄时，祂假装要继续前行；福音作者说：\BibleQuote{但祂自己假装要继续前行}，用了这个字眼，这正是说谎者极喜好的，以便他们可以无罪地说谎；仿佛一切假装都是谎言，而其实在真理的方式中，为了用一个事物表示另一个事物，许多事情通常是被假装的。如果耶稣假装要继续前行这件事，除此之外没有别的意义，那么有理由被判断为谎言；但若被正确理解并指向祂所愿意表示的事物，那就是一个奥秘。否则，凡因着某种要表示的事物的相似性，虽然从未发生过，却被叙述为发生过的一切，都将成为谎言。比如，关于一个人的两个儿子，与父亲同住的长子和去了远方的幼子，这个叙述得如此长的故事\BibleRef{路 15:11-32}。在这种虚构中，人们甚至将人的行为或言语赋予无理性的动物和无感觉的事物，以便通过这种假装的叙述而真实的意义，更迷人地暗示他们所愿之事。这不仅在世俗文学作者中，如\Person{贺拉斯}，老鼠对老鼠说话，黄鼠狼对狐狸说话，通过虚构的叙述而真实的意义指向所涉及的事；\Person{伊索}的类似寓言也指向同一目的，没有一个未受过教育的人会认为它们应被称为谎言；而且在圣经中，如\BibleBook{}{民长纪}，树木为自己寻找一位君王，对橄榄树、无花果树、葡萄树和荆棘说话\BibleRef{民 9:8-15}。这一切无疑是假装的，目的是通过假装的叙述（而非说谎的叙述）以真实的意义达到所想之事。我这样说是因为论及耶稣所记的：\BibleQuote{祂自己假装要继续前行}；免得任何人从这句话，像\Person{普里西利安主义者}那样，希望拥有说谎的许可，就争辩说除他人外连基督也说谎了。但若要明白祂假装所预表的，就请注意祂通过行动所实现的事。因为后来祂确实继续前行，升到诸天之上，却没有离弃祂的门徒。为了表示祂以后作为天主所行之事，此刻祂作为人假装如此行。因此，在假装前行中产生了一个真实的意义，因为在这次离去中，那意义的事实随后发生了。所以，否认祂通过行动实现了祂所表示的事实的人，才该争论基督通过假装说谎了。\par

29. 因这缘故，说谎的异端者在\OriginalTerm{新约圣经}{New Testament}中找不到任何适于模仿的说谎先例，他们便以为自己论点丰富，认为说谎是正当的；因为从旧约先知书中，由于除了少数明白人之外，看不出表意性的言语和行为（作为真实）应指向什么，他们便似乎为自己找出并提出了许多谎言。但为了拥有可用来辩护自己的、似乎适于模仿的欺骗先例，他们欺骗了自己，而\BibleQuote{他们的不义欺骗自己}。然而，那些在经中不被相信曾愿做先知的人，若在行为或言语中怀着欺骗的意愿假装了什么，即使可能从他们所行所说的事中，由预先放置其中并预定的全能者（祂知道如何甚至利用人的恶行来成就善）塑造出某种预言性的东西，但就这些人本身而言，无疑他们是说谎了。但他们不应仅仅因为出现在那些配称为圣的经书中就被视为适于模仿；因为这些书同时记载了人的恶行和善行：一些要避开，一些要效法；有些如此记载，以至于对其也下了定论；有些则没有明示判断，留给我们自行判断；因为我们不仅应被明显的事物滋养，也应被隐晦的事物操练。\par

30. 但这些人为甚么认为可以模仿说谎的\Person{塔玛尔}，却不认为可以模仿行奸淫的\Person{犹大}？因为他们都读过这两件事，而圣经对这两件事既未指责也未赞扬，只是叙述了二者，并将二者交由我们判断；但若它允许其中任何一件被模仿而免于惩罚，那才是惊人的。因为，我们知道，\Person{塔玛尔}说谎不是出于卖淫的欲望，而是出于怀胎的愿望。但是，奸淫，虽然\Person{犹大}的不是这样，但有些人的奸淫可能是为了使一个人得救，正如她的谎言是为了使一个人被怀胎；那么，是否可以因此行奸淫，如果因此认为说谎是正当的？因此，不仅关于说谎，而且关于人的一切行为，其中出现某种作为补偿的罪，我们必须考虑应当作出什么判断；免得我们不仅为任何小罪，而且为一切邪行敞开道路，以至于没有一种凶暴、可耻、渎神的行为，不能找出一个理由使它看似应该行事，从而使整个生活的正直被此观点颠覆。\par

31. 但若有人说有些谎言是正义的，就必须判断此人的话无异于说有些罪是正义的，从而有些正义之事也是不义的：还有什么比这更荒谬呢？因为一件事成为罪，岂不正是因为它违反正义吗？那么，当说有些罪大、有些罪小，因为这是真的；我们不要听信那些主张所有罪都相等的斯多葛派。但要说有些罪是不义的，有些是正义的，这岂不是等于说有不义的和正义的恶行？当宗徒\Person{若望}说：\BibleQuote{凡犯罪的，也行不义，罪就是不义。}因此，罪不可能是正义的，除非我们将罪之名用于另一件事，在那件事上人并不犯罪，而是为了罪而做或受某事。也就是说，为罪的祭献被称为\Quote{罪}，罪的惩罚有时也被称为罪。这些无疑可以被理解为正义的罪，当谈到正义的祭献或正义的惩罚时。但那些违反天主法律的事不能是正义的。有人对天主说：\BibleQuote{你的法律是真理。}因此，违反真理的事不能是正义的。现在有谁怀疑每个谎言都违反真理呢？所以不可能有正义的谎言。再者，有谁看不清楚每件正义的事都是出于真理呢？\Person{若望}呼喊说：\BibleQuote{没有谎言是出于真理。}所以没有谎言是正义的。因此，当\WorkTitle{圣经}向我们提出谎言的例子时，要么它们不是谎言，只是因未被理解而被当作谎言；要么，如果它们是谎言，也不值得效仿，因为它们不可能是正义的。\par

32. 但是，关于\WorkTitle{圣经}上所记的，天主善待了希伯来的接生婆和\Place{耶里哥}的妓女\Person{辣哈布}，这并不是因为她们说谎，而是因为她们怜悯了天主的子民。所以，在她们身上得到赏报的，不是她们的欺骗，而是她们的仁慈；是心灵的良善，而不是谎言之不义。因为，如果天主因她们后来所行的善工，愿意宽恕她们从前所犯的某些恶行，这并不奇怪和荒谬；同样，天主在同一时间、同一事上，同时看见了这两者，即出于怜悯的事和出于欺骗的事，便赏报了善，并为了这善而宽恕了那恶，这也不足为奇。因为如果那些出于肉欲而不是出于怜悯所犯的罪，因以后的怜悯善工而被赦免，那么那些出于怜悯本身所犯的罪，为何不能因怜悯的功劳而被赦免呢？因为以伤害为目的的罪，比以帮助为目的的罪更严重。因此，如果后者被随后的怜悯善工所涂抹，那么前者的罪较轻，为何不能因那人本身的怜悯而被涂抹呢？这怜悯既在他犯罪之前，又在他犯罪时伴随着他？这样看来似乎有理：但实际上，说\Quote{我不应该犯罪，但我要行怜悯的工作，以便能涂抹我以前所犯的罪，}是一回事；说\Quote{我应该犯罪，因为我无法以别的方式显示怜悯，}是另一回事。我说，说\Quote{我们已经犯了罪，让我们行善，}是一回事；说\Quote{让我们犯罪，好使善得以成就，}是另一回事。那里说的是\Quote{让我们行善，因为我们行了恶；}而这里说的是\Quote{让我们行恶，好使善得以产生。}因此，在那里我们要清除罪的污秽，在这里要警惕那教导人犯罪的道理。\par

33. 那么剩下的就是，我们要理解为那些妇女，无论是在\Place{埃及}还是在\Place{耶利哥}，她们因着人性和怜悯而获得了一个赏报，这赏报完全是暂时的，然而她们自己并不知道，这赏报却借着预言的象征预表了某种永恒的事物。但是，是否为了拯救人的性命而说谎，在任何情况下都是许可的，这个问题连最博学的人在解决时都感到疲倦，大大超出了那些可怜妇女的能力，她们置身于那些民族中间，并习惯了那些风俗。因此，天主耐心容忍了她们在这件事上的无知，也容忍了她们在其他同样不知道的事情上的无知——这些事情不是属于这个世界，而是属于未来世界的儿女所应当知道的。尽管如此，天主因着她们曾向祂的仆人所显示的人性仁慈，还是赐给了她们现世的赏报，虽然这些赏报也象征了一些天上的事物。而\Person{辣哈布}，从\Place{耶利哥}被救出后，进入了天主的子民之中，在那里，她不断进步，就可以获得永恒不朽的奖赏，这些奖赏是不能借任何谎言寻求的。然而，当她为以色列探子做了那件好事——就她的生活状况而言是值得称赞的工作——时，她还不是那样的人，以至于要求她\BibleQuote{在你的口中，是就是是，非就是非。}但是，至于那些接生婆，虽然她们是希伯来妇人，但如果她们只按肉体行事，那么她们从所得的现世赏报——就是天主为她们建立家室——中获得了什么或多大的好处呢？除非她们不断进步，到达了那座向天主歌唱的殿：\BibleQuote{住在你殿里的人是有福的，他们要永远赞美你。}然而，必须承认，它非常接近正义，虽然还不是现实，甚至现在就盼望和心性而言，那种心灵是值得称赞的，它从不撒谎，除非有意向和意愿对某人做好事，而不伤害任何人。但是，对于我们，当我们问一个好人有时是否可以说谎时，我们不是问关于一个属于\Place{埃及}、或\Place{耶利哥}、或\Place{巴比伦}、或仍属于地上的\Place{耶路撒冷}本身的人，这地上的\Place{耶路撒冷}与她的儿女同受奴役；而是问关于那座在上而自由的城的一位公民，那座城是我们的母亲，永在天上。对于我们的询问，回答是：\BibleQuote{谎言不是出于真理。}那城的儿子是真理的儿子。那城的儿子们是那些经上记载的人：\BibleQuote{在他们口中找不到谎言：}那城的儿子也是经上记载的人：\BibleQuote{接受言语的儿子必远离毁灭：但他接受，他就为自己接受了，从他口中没有虚假的话出来。}这些天上的\Place{耶路撒冷}和永恒圣城的儿子们，如果作为人，一旦任何谎言钻入他们心中，他们会谦卑地祈求宽恕，而不是进而从中寻求荣耀。\par

34. 但有人会说：那么那些希伯来产婆和辣哈布，若是拒绝说谎，不行怜悯，岂不是更好？绝非如此。那些希伯来妇女，如果她们是那种我们问是否应当说谎的人，便会避免说任何虚假的话，并且会最坦诚地拒绝杀害婴儿的卑贱服务。但你会说：她们自己会死。是的，但请看后果。她们会死，但为她们那无可比拟的更大赏报——天上的居所，远超她们在地上所建造的房屋——而死；她们会死，为最无罪的真理而忍受死亡后，进入永恒的福乐。那在耶里哥的妇人是怎样呢？她能这样做吗？如果她不说谎欺骗询问的市民，而是说实话出卖隐藏的客人，她岂不是会这样做？或者，她可以对那些询问者说：我知道他们在哪里，但我敬畏天主，我不会出卖他们。她确实可以说这话，如果她已经是一个没有诡诈的真正以色列女子；而她即将成为这样的人，藉着天主的仁慈，转往天主之城。但你会说，他们听到这，会杀死她，搜查房子。但他们也必定会找到她仔细隐藏的人吗？因为预见到这一点，那最谨慎的妇人已将他们放在一个地方，如果她说谎而不被相信，他们本可以不被发现。因此，即使她被同胞为怜悯之工所杀，她也会以在主眼中宝贵的死亡结束这必须结束的生命，并且她对他们的恩惠并非徒然。但你会说：\Quote{那些搜寻的人，若彻底搜查，来到她隐藏他们的地方，又该如何？} 以此方式可以说：如果一个最卑鄙无耻的妇人，不仅说谎，而且发假誓，却没有让他们相信她，又该如何？当然，即使如此，那些因恐惧而说谎的事仍有可能发生。我们又把天主的旨意和能力放在哪里？或许祂不能既保护她不对同胞说谎、不背叛天主的人，又保护祂的人免受一切伤害吗？因为藉着那在妇人说谎后保护他们的那一位，即使她没有说谎，祂也完全可以保护他们。除非我们忘记了这事在索多玛发生过：在那里，男人们怀着可憎的欲望向男人纵欲，却连他们寻找的人所在的房门都找不到；那时，那位义人，在极其相似的情形下，不愿为他的客人说谎——他不知他们是天使，却害怕他们遭受比死亡更甚的暴力。无疑，他本可以像耶里哥那妇人那样回答搜寻者。因为他们是同样地通过询问来搜寻。但那义人不愿意为了客人的身体而让自己的灵魂因说谎而玷污，却愿意为了这些身体而使自己女儿的身体被他人邪淫所侵犯。因此，人应当为人的暂时安全尽其所能；但当情况发展到除犯罪外无法保障他们的安全时，那么他就应认为自己已无能为力，因为他会发现只剩下他不该做的事。因此，关于耶里哥的辣哈布：她招待陌生人——天主的人；在招待他们时，她将自己置于危险中；她相信他们的天主；她勤谨地隐藏他们；她给予他们最忠实的劝告，让他们从另一条路回去——她应受赞美，堪为天上耶路撒冷的市民所效法。但在她说谎这件事上，虽然其中有些预言性的内容可被明智地解释，但并非作为可效法的榜样被智慧地提出；尽管天主荣耀地纪念了那些善事，却仁慈地忽略了这件恶事。\par

35. 既然事情如此，因为要详尽地处理狄克提纽斯那部\WorkTitle{Pound}中所列举的一切可效法的说谎先例未免太长，在我看来，以下是一条规则，不仅这些，而且所有此类事例都必须归到这条规则之下。即：要么，被认为说谎的事必须被证明并非如此；无论是当真理未被说出而虚假也未被说出时，还是当一种真实的示意让人理解另一事物时（这种比喻性的说法或做法在先知著作中比比皆是）。要么，那些被确证为谎言的，必须被证明并非可效法的；如果它们像其他罪过一样悄然潜入我们，我们不应将义德归于它们，而应请求宽恕。我确实如此认为，并且以上所辩论的内容也迫使我得出这一结论。\par

36. 但因为我们是人，生活在人中间，我承认自己尚未达到那些不受补偿罪困扰的人之列，所以在人间事上，我常常被人情所胜，无法抗拒别人对我说的话：\Quote{看，这里有一个病人，患重病危在旦夕，若将他唯一且最亲爱的儿子的死讯告诉他，他的体力将无法承受；他问你他的儿子是否还活着，而你知道他已经去世；你会如何回答？无论你说什么，除了以下三种之一：要么他死了，要么他还活着，要么我不知道；他只会相信他死了；因为他察觉到你害怕告诉他，又不愿说谎。}如果完全保持沉默，结果也一样。但那三种中，有两种是假的：他还活着，和我不知道；你只有说谎才能说出它们。而若你说出真实的那一种，即他死了，此人因此受到极大震动而死去，人们会喊叫说是你杀死了他。谁又能忍受人们指责他避免说谎可能救人一命，却选择了真理而害死人呢？我被这些反对意见深深打动，但也许并非明智。因为，当我将心思（姑且如此说）转向祂的理智之美——从祂的口中毫无虚假——尽管真理的光芒越发照亮我，我那软弱悸动的感官却退缩回去：但我被那无与伦美的爱火点燃，以至于轻视所有可能将我从那里召回的人世思虑。但这份爱能否坚持到在诱惑中不至于失去实效，却是一个大问题。当我默观那光明之善，其中绝无谎言之暗时，我不为以下说法所动：当我们拒绝说谎，而人因听到真理而死时，真理被称为凶手。因为，若一个荡妇向你求欢，你不答应，她因爱情的炽烈而死去，难道贞洁也成了凶手吗？或者，因为我们读到：\BibleQuote{我们在各处都是基督的馨香，在得救的人身上，也在丧亡的人身上}——对一些人，是生命的芬芳，对另一些人，是死亡的芬芳；难道我们也要称基督的馨香为凶手吗？但是，因为我们作为人，在这样的问题和矛盾中，大半被人情所胜或所困，所以宗徒随即又加上：\BibleQuote{这事谁能担当得起呢？}\par

37. 此外（这里更令人悲叹），如果我们一旦承认，为了挽救那个病人的生命而向他说谎说他的儿子还活着是正当的，那么，渐渐地、一步步地，邪恶就在我们身上增长，通过微小的积累，最终堆积成如此大量邪恶的谎言，在几乎难以察觉的侵蚀中，以至于再也找不到任何地方可以抵挡这巨大的灾祸，它因最小的积累而升涨成无限的力量。因此，经上极有远见地写道：\BibleQuote{轻视小事的人，必会渐渐跌倒。}不仅如此：那些如此迷恋此生，毫不犹豫地将它置于真理之上，为了使人不死——更确切地说，使一个终有一死的人晚一些死——的人，不仅要我们撒谎，甚至要我们发假誓！也就是说，免得人的虚妄健康稍纵即逝，我们竟要妄呼上主我们天主的名！有些博学之士甚至制定规则，划定界限，何时有义务发假誓，何时没有！哦，眼泪的源泉，你在哪里？我们该做什么？往哪里去？我们躲到哪里才能躲避真理的愤怒，如果我们不仅不回避谎言，还敢教导人发假誓？要知道，那些拥护和辩护谎言的人，要小心他们喜欢证明哪种或哪些谎言是正当的：至少在天主的敬礼中，让他们承认不可说谎；至少让他们远离发假誓和亵渎；至少在那里，在天主的圣名、天主作证、天主的誓言被用来起誓之处，在天主的宗教被谈论或对话之处，没有人说谎，没有人赞美、教导、命令、正当化谎言；至于其他种类的谎言，让那认为说谎正当的人自己挑选他认为最温和、最无罪的谎言。我知道，连教导说谎是合宜的人，也希望被认为是在教导真理。因为如果他教导的是假的，谁会在意听虚假的教义呢？在那教义中，教导者欺骗人，学习者也被欺骗。但是，为了找到门徒，他坚称自己是在教导真理，当他在教导说谎是合宜的时候，那个谎言怎能出自真理呢？因为若望宗徒宣告说：\BibleQuote{没有谎言是出于真理。}因此，有时说谎是正当的说法并不真实；而凡不真实的，绝不可说服任何人。\par

38. 但软弱为自己辩护，并借群众的赞许宣布自己拥有不可战胜的理由。于是它反驳说：\Quote{在人间，还有什么方法可以拯救那些陷于危险的人呢？这些人无疑会因为受骗而背离对他人或自身的致命伤害，如果我们作为人的情感不促使我们撒谎？}如果这群终有一死的、软弱的人群愿意耐心听我，我愿为真理说几句话。的确，至少虔诚、真诚、圣洁的贞洁无非是出于真理；凡是违背贞洁的，就是违背真理。那么，如果除此之外无法拯救陷于危险的人，我为何不去行淫呢？淫乱之所以违背真理，是因为它违背贞洁；然而，为了拯救陷于危险的人，我却说了一个最公开悖逆真理本身的谎言？贞洁在我们眼中如此值得尊重，而真理又何处得罪了我们？既然一切贞洁都出于真理，而且不仅是身体的贞洁，更是思想的贞洁就是真理，是的，连身体的贞洁也居于思想之中。最后，正如我刚刚说过并再说的：凡是为推荐和辩护任何谎言而反对我的人，他说的如果不是真理，又是什么？如果正因他说真理，所以应当被听，那他如何希望藉着说真理使我成为说谎者？谎言如何将真理当作自己的保护者？或者，真理为了征服自己的对手而取得胜利，以至自身被征服？谁能忍受这种荒谬？因此，我们绝不能说，那些断言有时可以撒谎的人，在断言这件事时是诚实的；否则（这是最荒谬、最愚蠢的信念），真理竟会教我们成为说谎者。这是何等的事：没有人从贞洁学到可以通奸；没有人从虔诚学到可以冒犯天主；没有人从仁慈学到可以伤害任何人；而为了说谎，我们却要从真理学习！但如果真理不教这个，它就不是真的；如果不是真的，就不值得学习；如果不值得学习，那么永远不值得说谎。\par

39. 但有人会说：\Quote{干粮是为成全的人。}因为许多事上，因着宽免，允许软弱有所松动，虽然就她（真理）的至诚而言，这些事毫不取悦于真理。让那不怕后果的人说这话吧；这后果是可怕的，一旦谎言以任何方式被允许。然而，绝不可允许它们攀升到如此高度，以致达到伪誓和亵渎；也不可提出任何理由，使犯伪誓成为合理，或更可憎的是，使亵渎天主成为合理。因为不能因为亵渎只是假装和谎言，就认为他没有被亵渎。因为照此说法，可以说犯伪誓不算犯罪，因为它是藉着谎言犯的：因为谁能藉着真理成为一个发假誓者？同样，没有人能藉着真理成为亵渎者。无疑，当一个人不知道所发誓的事情是假的，而相信它是真的时，这是一种较轻的假誓；就像\Person{扫禄}也因为无知而亵渎，更可原谅一样。但为什么亵渎比发假誓更严重呢？因为在假誓中，天主被召来见证一件虚假的事；但在亵渎中，虚假的事被直接说到天主自己。一个人，无论是发假誓者还是亵渎者，在多大程度上他们断言那些他们知道或相信是虚假的事情，就在多大程度上更不可原谅。因此，谁若说为了一个陷于危险的人现世的平安或生命可以说谎，若他还说为此甚至可以指天主起誓，甚至亵渎天主，他自己就太偏离永生的平安和生命的道路了。\par

40. 但有时他们将永恒救恩本身的危险抛向我们；他们叫嚷说，我们必须以谎言来抵御这危险，如果没有别的办法的话。例如，若一个将要领受圣洗的人落在不虔不敬、不信天主的恶人手中，若不欺骗他的看守者（以谎言）就无法接近他，使他获得重生之洗。面对这最令人反感的叫嚷——它迫使我们，不是为了人在这世上转瞬即逝的财富或尊荣，也不是为了现世的生命本身，而是为了一个人的永恒救恩，去说谎——我除了投奔于祢，真理啊，还能逃往何处？经由祢，贞洁呈现在我面前。因为，若这些看守者可以被诱骗（我们若行淫乱）而允许我们为那人施洗，我们为何拒绝行违反贞洁的事，却同意行违反真理的事（若以谎言就能欺骗他们）？毫无疑问，没有人会真诚地认为贞洁可爱，除非因为它是真理所命令的。那么，为了接近一个人为他施洗，就让看守者被谎言欺骗吧，如果真理这样命令。但真理怎能为了使人受洗而命令我们说谎呢？正如贞洁不命令我们为了使人受洗而犯奸淫。贞洁为何不这样命令？正是因为真理不这样教导。如果除了真理所教导的，我们就不该做，那么，当真理甚至不教导我们为了给人施洗而做出违反贞洁的事时，它又怎能教导我们为了给人施洗而做出违反它自身（真理）的事呢？但如同眼睛不够强壮不能直视太阳，却乐于观看被太阳照亮的事物，同样，已经有力以贞洁之美为乐的灵魂，却未必能直接在她自身——真理——中思虑那爱德从之得光的那一位，以致当他们要行某些违反真理的事时，他们不会像面对违反贞洁的事那样惊恐退缩。但那承受了圣言的儿子，他将远离丧亡，他口中不出虚假；他视被迫以谎言救助邻人，与被迫以淫行救助同样不可为。圣父垂听并允准他的祈祷，使他能够无需谎言就救助圣父自己（其判断不可测度）愿意救助的人。因此，这样的儿子防备谎言，如同防备罪恶。因为有时谎言之名被用作罪恶之名；经上说：\BibleQuote{所有的人都是撒谎者。}这话说得就好像是说：\BibleQuote{所有的人都是罪人。}还有：\BibleQuote{但若天主的真实因我的谎言而越发丰盛。}因此，当人如常人一样说谎时，他如常人一样犯罪，并要受那判语：\BibleQuote{所有的人都是撒谎者}以及\BibleQuote{若我们说我们没有罪，便是自欺，真理不在我们内。}但当他口中不出虚假时，那是按那恩宠成就的，关于这恩宠经上说：\BibleQuote{凡由天主生的，就不犯罪。}因为若这出生单独在我们内，就没有人会犯罪；当它单独存在时，也没有人会犯罪。但现在，我们仍拖带着那所生的可朽坏的身体；虽然，按我们所重生的，若我们正直行走，我们日日在内里更新。但当这可朽坏穿上不可朽坏时，生命将完全吞没它，不再有死亡的刺存留。这死亡的刺就是罪。\par

41. 那么，我们要么借着行义来避免说谎，要么借着悔改来承认它们；但绝不能在我们生活中不幸地充斥着谎言的同时，还通过教导使它们增多。不过，让那持这种想法的人，选择他愿意的哪种谎言，以便他能帮助处于危险中的邻人达到他所期望的救恩；但即使对这样的人，我们也要求，在任何情况下都不可被引向虚假宣誓和亵渎。至少让我们判断这些恶行要么比淫乱行为更严重，要么肯定不亚于它们。确实值得思考的是，人们常常在怀疑妻子犯奸淫时，要求她们起誓；他们这样做当然是因为相信，即使那些不惧犯奸淫的人，也可能害怕伪证。事实上，有些淫荡的妇女，虽然不怕以非法拥抱欺骗丈夫，却害怕向那些被她们欺骗的丈夫虚假地呼求天主作证。那么，一个贞洁而虔诚的人，怎么会不愿意通过奸淫来帮助一个人领受圣洗，却愿意通过伪证来帮助他——而伪证连犯奸淫者通常都害怕呢？此外，如果通过伪证这样做是可憎的，那么通过亵渎岂不是更可憎？因此，一个基督徒绝不能否认和亵渎基督，以便使另一个人成为基督徒；并且通过丧失自己来寻求一个，如果他教导这样的人这类事情，可能会使找到的人反而丧亡。那么，那本被称为\WorkTitle{镑}的书，你必须用这种方法来驳斥和摧毁；即它的那个标题中，他们教条地主张可以为隐瞒宗教而说谎，你必须明白这必须首先被切除；这样一来，他们为努力推崇圣经作为他们谎言的赞助者而使用的那些证言，你必须证明一部分并不是谎言，一部分即使是谎言，也不值得模仿；并且，如果软弱为自己擅自如此，以至于有些真理不认可的事被可宽恕地允许给她，但你必须坚定不移地持守并捍卫：在神圣的宗教中，绝无任何时间允许说谎。至于隐藏的异端者，正如我们不会通过犯奸淫来发现隐藏的奸淫者，也不会通过谋杀来发现谋杀者，也不会通过行巫术来发现行巫术者，同样，我们也不应通过说谎来发现说谎者，或通过亵渎来发现亵渎者：根据我们在本书中如此详尽地陈述的推理，我们终于艰难地达到了我们在此处设定的目标。\par

\SourceRecord{Translated by H. Browne.；Nicene and Post-Nicene Fathers, First Series；Vol. 3.；Edited by Philip Schaff.；Buffalo, NY: Christian Literature Publishing Co.,；1887.；<http://www.newadvent.org/fathers/1313.htm>.}
